% GENERATED FILE. Edit canonical lessons, then run npm run generate:books.
% canonical-source-hash: fnv1a64:547be9b00963c27b
% canonical-lessons: ML-C547-pailarru, ML-C547-dantavaidyan, ML-C547-vastusilpi, ML-C547-ilaktrisyan, ML-C547-plambar

\chapter{Pilots, Plumbers and Scientists}
\label{ch:ml-pilots-plumbers-and-scientists}
\hlchaptermodality{547}

\begin{chapteropening}
\textbf{By the end of this chapter:} \emph{I can say a pilot, a dentist, an architect, an electrician and a plumber.}

Complete the last lesson of chapter 547: I can say a pilot, a dentist, an architect, an electrician and a plumber.
\end{chapteropening}

\section[\texorpdfstring{\ml{പൈലറ്റ്} (pailaṟṟŭ)}{പൈലറ്റ് (pailaṟṟŭ)}]{\ml{പൈലറ്റ്} (pailaṟṟŭ) — a pilot}

\label{lesson:ML-C547-pailarru}



\begin{quote}
Before the new one: say the Malayalam for work experience, then the Malayalam for to lead.
\end{quote}

\subsection*{You'll want to know: \ml{പൈലറ്റ്}}
\textbf{\ml{പൈലറ്റ്}} — \emph{pailaṟṟŭ} — \textquotedblleft{}a pilot\textquotedblright{}.

\begin{quote}
\ml{കുട്ടിക്കാലം} \ml{മുതൽ} \ml{അവന്} \ml{പൈലറ്റ്} \ml{ആകാനായിരുന്നു} \ml{ആഗ്രഹം}.
\end{quote}

\begin{quote}
\emph{kuṭṭikkālaṁ mutal avanŭ pailaṟṟŭ ākānāyirunnu āgrahaṁ.}
\end{quote}

\begin{quote}
Since childhood he had wanted to be a pilot.
\end{quote}

To fly a plane is \ml{വിമാനം} \ml{പറത്തുക}.

\begin{cousinweb}
\textbf{\ml{പൈലറ്റ്}} (\emph{pailaṟṟŭ}) was borrowed from English \emph{pilot}.
\end{cousinweb}

\begin{grammarlens}[title={What you've built}]
One more word for describing what you lived and felt.
\end{grammarlens}

\subsection*{Your turn}
\begin{itemize}
\raggedright
  \item \emph{Say it:} \emph{pailaṟṟŭ}
  \item \emph{Say it:} \emph{pailaṟṟŭ}, once more
  \item \emph{Say it:} \emph{kuṭṭikkālaṁ mutal avanŭ pailaṟṟŭ ākānāyirunnu āgrahaṁ.} — Since childhood he had wanted to be a pilot.
  \item \emph{Recall:} say the Malayalam for work experience, then the Malayalam for to lead, then say \emph{pailaṟṟŭ} again
\end{itemize}

\begin{tcolorbox}[breakable,colback=teal!4,colframe=teal!35!black,title={Before you move on}]
What does \ml{പൈലറ്റ്} mean? (\textquotedblleft{}A pilot\textquotedblright{}.) Say it once more.
\end{tcolorbox}
\section[\texorpdfstring{\ml{ദന്തവൈദ്യൻ} (dantavaidyan)}{ദന്തവൈദ്യൻ (dantavaidyan)}]{\ml{ദന്തവൈദ്യൻ} (dantavaidyan) — a dentist}

\label{lesson:ML-C547-dantavaidyan}



\begin{quote}
Before the new one: say the Malayalam for to lead, then the Malayalam for a pilot.
\end{quote}

\subsection*{You'll want to know: \ml{ദന്തവൈദ്യൻ}}
\textbf{\ml{ദന്തവൈദ്യൻ}} — \emph{dantavaidyan} — \textquotedblleft{}a dentist\textquotedblright{}.

\begin{quote}
\ml{പല്ലുവേദന} \ml{കാരണം} \ml{ഞാൻ} \ml{ദന്തവൈദ്യനെ} \ml{കണ്ടു}.
\end{quote}

\begin{quote}
\emph{palluvēdana kāraṇaṁ ñān dantavaidyane kaṇṭu.}
\end{quote}

\begin{quote}
I saw the dentist because of a toothache.
\end{quote}

\ml{ദന്തം} a tooth + \ml{വൈദ്യൻ} a doctor.

\begin{cousinweb}
\textbf{\ml{ദന്തവൈദ്യൻ}} (\emph{dantavaidyan}) is a learned compound built from Sanskrit \emph{danta}, \textquotedblleft{}a tooth\textquotedblright{} (CDIAL 6152), and \emph{vaidya}, \textquotedblleft{}a physician\textquotedblright{} (CDIAL 12140).
\end{cousinweb}

\begin{grammarlens}[title={What you've built}]
One more word for describing what you lived and felt.
\end{grammarlens}

\subsection*{Your turn}
\begin{itemize}
\raggedright
  \item \emph{Say it:} \emph{dantavaidyan}
  \item \emph{Say it:} \emph{dantavaidyan}, once more
  \item \emph{Say it:} \emph{palluvēdana kāraṇaṁ ñān dantavaidyane kaṇṭu.} — I saw the dentist because of a toothache.
  \item \emph{Recall:} say the Malayalam for to lead, then the Malayalam for a pilot, then say \emph{dantavaidyan} again
\end{itemize}

\begin{tcolorbox}[breakable,colback=teal!4,colframe=teal!35!black,title={Before you move on}]
What does \ml{ദന്തവൈദ്യൻ} mean? (\textquotedblleft{}A dentist\textquotedblright{}.) Say it once more.
\end{tcolorbox}
\section[\texorpdfstring{\ml{വാസ്തുശില്പി} (vāstuśilpi)}{വാസ്തുശില്പി (vāstuśilpi)}]{\ml{വാസ്തുശില്പി} (vāstuśilpi) — an architect}

\label{lesson:ML-C547-vastusilpi}



\begin{quote}
Before the new one: say the Malayalam for a pilot, then the Malayalam for a dentist.
\end{quote}

\subsection*{You'll want to know: \ml{വാസ്തുശില്പി}}
\textbf{\ml{വാസ്തുശില്പി}} — \emph{vāstuśilpi} — \textquotedblleft{}an architect\textquotedblright{}.

\begin{quote}
\ml{ഞങ്ങളുടെ} \ml{പുതിയ} \ml{വീടിന്റെ} \ml{പ്ലാൻ} \ml{ഒരു} \ml{വാസ്തുശില്പി} \ml{തയ്യാറാക്കി}.
\end{quote}

\begin{quote}
\emph{ñaṅṅaḷuṭe putiya vīṭinṟe plān oru vāstuśilpi tayyāṟākki.}
\end{quote}

\begin{quote}
An architect drew up the plan of our new house.
\end{quote}

\ml{വാസ്തു}, a building site, + the Sanskrit word for a craftsman: the one who designs the building.

\begin{cousinweb}
\textbf{\ml{വാസ്തുശില്പി}} (\emph{vāstuśilpi}) is a learned compound built from Sanskrit \emph{vāstu}, \textquotedblleft{}the site of a house\textquotedblright{} (CDIAL 11606), and \emph{śilpin}, \textquotedblleft{}skilled in art, an artisan\textquotedblright{} (CDIAL 12471).
\end{cousinweb}

\begin{grammarlens}[title={What you've built}]
One more word for describing what you lived and felt.
\end{grammarlens}

\subsection*{Your turn}
\begin{itemize}
\raggedright
  \item \emph{Say it:} \emph{vāstuśilpi}
  \item \emph{Say it:} \emph{vāstuśilpi}, once more
  \item \emph{Say it:} \emph{ñaṅṅaḷuṭe putiya vīṭinṟe plān oru vāstuśilpi tayyāṟākki.} — An architect drew up the plan of our new house.
  \item \emph{Recall:} say the Malayalam for a pilot, then the Malayalam for a dentist, then say \emph{vāstuśilpi} again
\end{itemize}

\begin{tcolorbox}[breakable,colback=teal!4,colframe=teal!35!black,title={Before you move on}]
What does \ml{വാസ്തുശില്പി} mean? (\textquotedblleft{}An architect\textquotedblright{}.) Say it once more.
\end{tcolorbox}
\section[\texorpdfstring{\ml{ഇലക്ട്രീഷ്യൻ} (ilakṭrīṣyan)}{ഇലക്ട്രീഷ്യൻ (ilakṭrīṣyan)}]{\ml{ഇലക്ട്രീഷ്യൻ} (ilakṭrīṣyan) — an electrician}

\label{lesson:ML-C547-ilaktrisyan}



\begin{quote}
Before the new one: say the Malayalam for a dentist, then the Malayalam for an architect.
\end{quote}

\subsection*{You'll want to know: \ml{ഇലക്ട്രീഷ്യൻ}}
\textbf{\ml{ഇലക്ട്രീഷ്യൻ}} — \emph{ilakṭrīṣyan} — \textquotedblleft{}an electrician\textquotedblright{}.

\begin{quote}
\ml{ഇലക്ട്രീഷ്യൻ} \ml{വന്ന്} \ml{ഫാൻ} \ml{നന്നാക്കി}.
\end{quote}

\begin{quote}
\emph{ilakṭrīṣyan vannŭ phān nannākki.}
\end{quote}

\begin{quote}
The electrician came and fixed the fan.
\end{quote}

English -ician comes out as -\ml{ീഷ്യൻ} in Malayalam: \ml{ഇലക്ട്രീഷ്യൻ}, electrician.

\begin{cousinweb}
\textbf{\ml{ഇലക്ട്രീഷ്യൻ}} (\emph{ilakṭrīṣyan}) was borrowed from English \emph{electrician}.
\end{cousinweb}

\begin{grammarlens}[title={What you've built}]
One more word for describing what you lived and felt.
\end{grammarlens}

\subsection*{Your turn}
\begin{itemize}
\raggedright
  \item \emph{Say it:} \emph{ilakṭrīṣyan}
  \item \emph{Say it:} \emph{ilakṭrīṣyan}, once more
  \item \emph{Say it:} \emph{ilakṭrīṣyan vannŭ phān nannākki.} — The electrician came and fixed the fan.
  \item \emph{Recall:} say the Malayalam for a dentist, then the Malayalam for an architect, then say \emph{ilakṭrīṣyan} again
\end{itemize}

\begin{tcolorbox}[breakable,colback=teal!4,colframe=teal!35!black,title={Before you move on}]
What does \ml{ഇലക്ട്രീഷ്യൻ} mean? (\textquotedblleft{}An electrician\textquotedblright{}.) Say it once more.
\end{tcolorbox}
\section[\texorpdfstring{\ml{പ്ലംബർ} (plambar)}{പ്ലംബർ (plambar)}]{\ml{പ്ലംബർ} (plambar) — a plumber}

\label{lesson:ML-C547-plambar}



\begin{quote}
Before the new one: say the Malayalam for an architect, then the Malayalam for an electrician.
\end{quote}

\subsection*{You'll want to know: \ml{പ്ലംബർ}}
\textbf{\ml{പ്ലംബർ}} — \emph{plambar} — \textquotedblleft{}a plumber\textquotedblright{}.

\begin{quote}
\ml{പൈപ്പ്} \ml{ചോർന്നപ്പോൾ} \ml{ഞങ്ങൾ} \ml{പ്ലംബറെ} \ml{വിളിച്ചു}.
\end{quote}

\begin{quote}
\emph{paippŭ cōrnnappōḷ ñaṅṅaḷ plambaṟe viḷiccu.}
\end{quote}

\begin{quote}
When the pipe leaked, we called the plumber.
\end{quote}

The object form adds -\ml{എ}: \ml{പ്ലംബറെ} \ml{വിളിച്ചു}.

\begin{cousinweb}
\textbf{\ml{പ്ലംബർ}} (\emph{plambar}) was borrowed from English \emph{plumber}.
\end{cousinweb}

\begin{grammarlens}[title={What you've built}]
One more word for describing what you lived and felt.
\end{grammarlens}

\subsection*{Your turn}
\begin{itemize}
\raggedright
  \item \emph{Say it:} \emph{plambar}
  \item \emph{Say it:} \emph{plambar}, once more
  \item \emph{Say it:} \emph{paippŭ cōrnnappōḷ ñaṅṅaḷ plambaṟe viḷiccu.} — When the pipe leaked, we called the plumber.
  \item \emph{Recall:} say the Malayalam for an architect, then the Malayalam for an electrician, then say \emph{plambar} again
\end{itemize}

\begin{tcolorbox}[breakable,colback=teal!4,colframe=teal!35!black,title={Before you move on}]
What does \ml{പ്ലംബർ} mean? (\textquotedblleft{}A plumber\textquotedblright{}.) Say it once more.
\end{tcolorbox}
